\linespread{1.1}         			% Palatino needs more leading (space between lines)
\chapter*{Conclusion et Perspectives}
%%%%%%%%%%%%%%%%%%%%%%%%%%%%%%%%%%%%
\label{Chapitre-8}
\linespread{1.5}         			% Palatino needs more leading (space between lines)
%	\minitoc
%	\newpage

\addcontentsline{toc}{chapter}{Conclusion et Perspectives}


Dans ce dernier chapitre, nous allons faire un dernier bilan en effectuant une synth�se des �tapes r�alis�es durant nos travaux. Ensuite, nous d�gagerons les perspectives et les �volutions futures concernant ses travaux.

N'oubliez pas que l'introduction et la conclusion doivent se r�pondre l'une � l'autre. L'introduction pose les points que vous allez aborder, les questions auxquelles vous allez r�pondre. La conclusion, elle, donne une synth�se, dresse un bilan. Une fois que vous aurez r�dig� les 2, il faudra les relire � la suite pour v�rifier qu'il y a coh�rence.

\noindent La conclusion comprend 3 �l�ments importants :

\begin{enumerate}
	\item Rappel de l'id�e directrice : Reprendre l'id�e directrice du travail � r�aliser, en comman�ant par rappeler le contexte et la probl�matique.
	\item Faire une synth�se des travaux r�alis�s, en mettant le point sur les r�sultats obtenus.
	\item Et terminer par les perspectives en �vocant des pistes de recherches futures.\\
\end{enumerate}






